\documentclass[11pt, a4paper]{article}

% --- PACCHETTI ---
\usepackage[utf8]{inputenc}
\usepackage[italian]{babel}
\usepackage[margin=2.5cm]{geometry} % Imposta i margini della pagina
\usepackage{graphicx}               % Per inserire il logo
\usepackage{tabularx}               % Per tabelle larghe quanto la pagina
\usepackage[table]{xcolor}          % Per colorare le celle delle tabelle
\usepackage[hidelinks]{hyperref}    % Per stile link cliccabili
\usepackage{lmodern}      % Rende i font scalabili all'infinito
\usepackage{anyfontsize}  % Permette di usare dimensioni personalizzate

\begin{document}

% ==============================================================================
% FRONTESPIZIO
% ==============================================================================
\begin{titlepage}
    \centering
    \vspace*{1cm}

    % LOGHI
    \includegraphics[width=0.45\textwidth]{../../../../assets/logo_unipd.jpg}
    \hfill
    \includegraphics[width=0.45\textwidth]{../../../../assets/logo.png}\par
    \vspace{0.5cm}

    % NOME DEL GRUPPO
    {\fontsize{45}{50}\sffamily\selectfont \textbf{On-By-One}\par} 
    \vspace{1.5cm}

    % TITOLO DEL VERBALE
    {\huge \bfseries Verbale interno [YYYY/MM/DD] \par}
    \vspace{2cm}

    % MEMBRI DEL GRUPPO
    {\large
    \textbf{On-By-One}:\par
    \vspace{0.3cm}
    Sara Ballardin, Shuhui Chen, Simone Giuliano, Alvise Greggio,
    
    Daniele Nuzzaci, Federico Temporin, Tommaso Zorzetto 
    
    (2101048, 2138001, 2145591, 2141004, 2139039, 2146764, 2148629)
    \par}

    \vfill

    % TABELLA INFORMAZIONI DOCUMENTO 
    \begin{table}[h]
        \centering
        \renewcommand{\arraystretch}{1.5}
        \begin{tabularx}{\textwidth}{|c|c|c|X|}
            \hline
            \rowcolor{gray!20} \multicolumn{4}{|c|}{\textbf{Informazioni documento}} \\
            \hline
            \rowcolor{gray!20} \textbf{Versione} & \textbf{Data} & \textbf{Stato} & \multicolumn{1}{c|}{\textbf{Destinatari}} \\
            \hline
            [X.X.X] & [YYYY/MM/DD] & [Bozza/Approvato] & On-By-One, prof. Tullio Vardanega, prof. Riccardo Cardin \\
            \hline
        \end{tabularx}
    \end{table}

    \vspace{1cm}
    
    % CONTATTI
    {\large Contatti: \href{mailto:onbyone.swe@gmail.com}{onbyone.swe@gmail.com} }
    
\end{titlepage}

\newpage

% ==============================================================================
% INDICE
% ==============================================================================
\tableofcontents
\newpage

% ==============================================================================
% INFORMAZIONI GENERALI
% ==============================================================================
\section{Informazioni generali}

    \begin{table}[h]
    \centering
    \renewcommand{\arraystretch}{1.5}
    \begin{tabularx}{\textwidth}{|c|>{\centering\arraybackslash}X|}
        \hline
        \rowcolor{gray!20} \multicolumn{2}{|c|}{\textbf{Dettagli della riunione}} \\
        \hline
        \textbf{Data} & [YYYY/MM/DD] \\ \hline
        \textbf{Ora inizio} & [HH:MM] \\ \hline
        \textbf{Ora fine} & [HH:MM] \\ \hline
        \textbf{Modalità} & [Discord / In presenza] \\ \hline
    \end{tabularx}
\end{table}

\begin{table}[h]
    \centering
    \renewcommand{\arraystretch}{1.5}
    \begin{tabularx}{\textwidth}{|>{\centering\arraybackslash}X|>{\centering\arraybackslash}X|>{\centering\arraybackslash}X|}
        \hline
        \rowcolor{gray!20} \multicolumn{3}{|c|}{\textbf{Partecipanti}} \\
        \hline
        \rowcolor{gray!20} \textbf{Nome e cognome} & \textbf{Ruolo} & \textbf{Presenza} \\
        \hline
        Sara Ballardin & [Ruolo] & Presente \\ \hline
        Shuhui Chen & [Ruolo] & Presente \\ \hline
        Simone Giuliano & [Ruolo] & Presente \\ \hline
        Alvise Greggio & [Ruolo] & Presente \\ \hline
        Daniele Nuzzaci & [Ruolo] & Presente \\ \hline
        Federico Temporin & [Ruolo] & Presente \\ \hline
        Tommaso Zorzetto & [Ruolo] & Presente \\ \hline
    \end{tabularx}
\end{table}

% ==============================================================================
% CORPO DEL VERBALE
% ==============================================================================

\section{Ordine del giorno}
\begin{itemize}
    \item Punto 1 dell'ordine del giorno
    \item Punto 2 dell'ordine del giorno
    \item Punto 3 dell'ordine del giorno
\end{itemize}

\section{Discussioni}
[Riassunto delle discussioni affrontate per ogni punto dell'ordine del giorno.]

\section{Decisioni}
\begin{itemize}
    \item \textbf{[VI-YYYY-MM-DD.1]} [Decisione 1.]
    \item \textbf{[VI-YYYY-MM-DD.2]} [Decisione 2.]
\end{itemize}

\section{To-do}
\begin{center}
    \renewcommand{\arraystretch}{1.5}
    \begin{tabularx}{\textwidth}{|c|X|c|c|}
        \hline
        \rowcolor{gray!20} 
        \multicolumn{1}{|c|}{\textbf{ID issue}} & 
        \multicolumn{1}{c|}{\textbf{Task}} & 
        \multicolumn{1}{c|}{\textbf{Assegnatario}} & 
        \multicolumn{1}{c|}{\textbf{Scadenza}} \\
        \hline
        [- o \#numero] & [Descrizione] & [Nome e cognome] & [YYYY/MM/DD] \\
        \hline
        [- o \#numero] & [Descrizione] & [Nome e cognome] & [YYYY/MM/DD] \\
        \hline
    \end{tabularx}
\end{center}

% ==============================================================================
% REGISTRO DELLE MODIFICHE
% ==============================================================================
\section{Registro delle modifiche}
\begin{center}
    \renewcommand{\arraystretch}{1.5}
    \begin{tabularx}{\textwidth}{|c|c|c|X|c|}
        \hline
        \rowcolor{gray!20} 
        \multicolumn{1}{|c|}{\textbf{Versione}} & 
        \multicolumn{1}{c|}{\textbf{Data}} & 
        \multicolumn{1}{c|}{\textbf{Autore}} & 
        \multicolumn{1}{c|}{\textbf{Descrizione}} & 
        \multicolumn{1}{c|}{\textbf{Verificatore}} \\
        \hline
        [1.X.X] & [YYYY/MM/DD] & [Nome e cognome] & [Descrizione] & [Nome e cognome] \\
        \hline
        [0.1.0] & [YYYY/MM/DD] & [Nome e cognome] & [Descrizione] & [Nome e cognome] \\
        \hline
        
    \end{tabularx}
\end{center}

\end{document}
